\documentclass[UTF8,aspectratio=169,11pt]{ctexbeamer}
\usepackage{amsmath,amssymb}
\usepackage{graphicx}
\usepackage{booktabs}
\usepackage{tabularx}
\usepackage{array}
\usepackage{xcolor}
\usepackage{hyperref}
\definecolor{primary}{HTML}{17365D}
\definecolor{accent}{HTML}{1F77B4}
\definecolor{evidence}{HTML}{1E8449}
\definecolor{warning}{HTML}{C55A11}
\definecolor{muted}{HTML}{5B6573}
\setbeamersize{text margin left=0.55cm,text margin right=0.55cm}
\setlength{\emergencystretch}{2em}
\setbeamertemplate{navigation symbols}{}{}
\setbeamercolor{normal text}{fg=black!86,bg=white}
\setbeamercolor{structure}{fg=primary}
\setbeamercolor{frametitle}{fg=primary,bg=white}
\setbeamercolor{block title}{fg=primary,bg=primary!6}
\setbeamercolor{block body}{fg=black!86,bg=primary!3}
\setbeamerfont{frametitle}{size=\large,series=\bfseries}
\setbeamertemplate{footline}{%
  \hbox{\begin{beamercolorbox}[wd=\dimexpr\paperwidth-1.1cm\relax,ht=2.3ex,dp=0.9ex,
    leftskip=0.55cm,rightskip=0.55cm]{author in head/foot}%
    \scriptsize\color{muted}PyQCD 13-conf TMD production\hfill
    \insertframenumber/\inserttotalframenumber
  \end{beamercolorbox}}}
\newcommand{\source}[1]{\vspace{0.06em}\par{\raggedright\scriptsize\color{muted}来源：#1\par}}
\title[PyQCD TMD]{梯度流重整化核子胶子 TMD-PDF 全链计算}
\author[PyQCD]{PyQCD production run}
\date{2026-11-01}
\begin{document}
\begin{frame}[plain]{梯度流重整化核子胶子 TMD-PDF：13 组态全链完成，但物理接受仍未闭合}
  \centering\begin{minipage}{0.86\linewidth}\centering{\Large\bfseries PyQCD 真实数据计算结果}\par\vspace{0.7em}13 configurations \quad $\tau/a^2=1,2,3$ \quad $z/a=0\ldots12$ \quad $b/a=0\ldots6$\par\vspace{1em}\textcolor{warning}{状态：diagnostic\_only；physical\_acceptance=false}\end{minipage}

\end{frame}

\begin{frame}[t]{结论先行：格点数据链已产出，但只能作为诊断性 TMD-PDF}
  \begin{columns}[T,totalwidth=\textwidth]
\begin{column}{0.32\textwidth}\begin{block}{已完成}12 个 1D ratio、6 个 2D ratio、6 个 PDF/TMD-PDF 产出。\end{block}\end{column}
\begin{column}{0.32\textwidth}\begin{block}{最强证据}13/13 组态、859 个数据文件；总字节 57914374503。\end{block}\end{column}
\begin{column}{0.32\textwidth}\begin{block}{未闭合}soft/rapidity、微扰 matching、胶子--夸克混合、连续极限均未完成。\end{block}\end{column}\end{columns}

  \source{汇总：\nolinkurl{/root/PyQCD/data/runs/tmd-ext-20261103/analysis/analysis_summary.json}}
\end{frame}

\begin{frame}[t]{验收标准区分“数据链完整”与“物理方案闭合”}
  \small\begin{tabularx}{\linewidth}{@{}l X l@{}}\toprule 层级 & 本次证据 & 状态 \\ \midrule 输入与缓存 & 13 组原始 gauge/eigen/peram 与 canonical HDF5 全链 & 确证 \\ 统计与图表 & 1D/2D ratio、JK、平台、Fourier、CS 与误差账本 & 确证 \\ soft/rapidity & 无同几何测量 & 未验证 \\ 微扰 matching & 按要求跳过非格点 matching & 不适用 \\ 连续极限 & 单格距、单系综 & 未验证 \\ \bottomrule\end{tabularx}

  \source{验收定义来自 pyqcd-tmd-algorithm validation reference。}
\end{frame}

\begin{frame}[t]{全链以同一组态索引贯穿流化、算符、统计和 Fourier}
  \small\begin{tabularx}{\linewidth}{@{}l X@{}}
\toprule 阶段 & 输入 $\to$ 输出 \\
\midrule
蒸馏顶点 & eigenvector $\to$ VdV/VVV canonical HDF5 \\
外态关联 & perambulator + VVV $\to$ source-resolved C2 与连接 3pt \\
Wilson flow & raw gauge $\to$ $V_\tau$、$E(\tau)$ 与 $t_0$ \\
胶子算符 & flowed gauge $\to$ 三方向 1D/2D OPE \\
比值与统计 & C2 + OPE $\to$ disconnected ratio、JK 样本、c0/hR \\
格点提取 & hR $\to$ 准 PDF、准 TMD-PDF、两动量 CS 核 \\
\bottomrule\end{tabularx}

  \source{\texttt{\detokenize{pyqcd/testing/tmd/run_full_production.py:617}}}
\end{frame}

\begin{frame}[t]{输入严格绑定 13 个目标组态、三流时和三方向位移网格}
  \small\begin{tabularx}{\linewidth}{@{}l X@{}}\toprule 参数 & 设置 \\ \midrule 组态 1 & 4150, 6250, 6450, 6650, 6850, 7050, 7250 \\ 
组态 2 & 7450, 7650, 7850, 8050, 8250, 8450 \\ 
格点 & $24^3\times72$, $a=0.1053$ fm \\ 动量 & P000, P200, P020, P002, P400 \\ 流时/步长 & $\tau/a^2=1,2,3$, RK3 $\epsilon=0.01$ \\ 非局域几何 & $z=0\ldots12$, $b=0\ldots6$, $L=12a$ \\ 二维几何 & $b_x=0\ldots6$, $b_y=0\ldots2$, 两种 staircase \\ \bottomrule\end{tabularx}

  \source{运行参数写入 workers/*.json 与 canonical cache contract。}
\end{frame}

\begin{frame}[t]{4 个 worker 分片执行；CPU/RAM/GPU 显存配置门为 80\%}
  \small\begin{tabularx}{\linewidth}{@{}>{\raggedright\arraybackslash}p{0.20\linewidth} l l >{\raggedright\arraybackslash}X l@{}}
\toprule
worker & backend/device & 状态 & 组态 & 完成时间 \\
\midrule
worker\_ext-a & torch/0 & complete & 4150 & 2026-11-03T18:15:09 \\
worker\_p100-a & cupy/0 & complete & 6250, 6850, 7450, 8050 & 2026-11-02T07:53:56 \\
worker\_p100-b & cupy/0 & complete & 6450, 7050, 7650, 8250 & 2026-11-02T06:15:52 \\
worker\_torch\_0 & torch/0 & complete & 6650, 7250, 7850, 8450 & 2026-11-02T11:20:54 \\
\bottomrule
\end{tabularx}\vspace{0.5em}
\begin{tabularx}{0.92\linewidth}{@{}c X c c c@{}}
\toprule
GPU & 设备 & 显存峰值 & 平均利用率 & 瞬时峰值 \\
\midrule
0 & worker device 0 & 32.07\% & 不可用 & 不可用 \\
1 & worker device 1 & 7.305\% & 不可用 & 不可用 \\
2 & worker device 2 & 7.305\% & 不可用 & 不可用 \\
\bottomrule
\end{tabularx}

  \source{CPU峰值 36.6\%，RAM峰值 22.1\%，GPU显存峰值 32.07\%，磁盘峰值 86.35\%。worker duty-cycle 限制平均负载。系统级 GPU 遥测瞬时峰值 100\%，因此不能认证全局瞬时利用率始终不超过 80\%。}
\end{frame}

\begin{frame}[t]{Wilson flow 显示平稳 UV 平滑并给出 $t_0$ 诊断}
  \includegraphics[width=0.88\linewidth,height=0.62\textheight,keepaspectratio]{/root/PyQCD/data/runs/tmd-ext-20261103/analysis/flow/flow_history.png}

\source{13 组态 Wilson-flow 作用量密度与 \$t\textasciicircum{}2\textbackslash{}langle E\textbackslash{}rangle\$；数据文件：flow\_history.png}

  \source{\texttt{\detokenize{pyqcd/renorm/_gradient_flow.py:272}}}
\end{frame}

\begin{frame}[t]{核子 2pt 的有效质量与平台窗作为诊断量报告}
  \begin{columns}[T,totalwidth=\textwidth]
\begin{column}{0.49\textwidth}\includegraphics[width=0.96\linewidth,height=0.62\textheight,keepaspectratio]{/root/PyQCD/data/runs/tmd-ext-20261103/analysis/spectrum/effective_mass_nopol_P000.png}

\source{P000 nopol effective mass；数据文件：effective\_mass\_nopol\_P000.png}\end{column}
\begin{column}{0.49\textwidth}\includegraphics[width=0.96\linewidth,height=0.62\textheight,keepaspectratio]{/root/PyQCD/data/runs/tmd-ext-20261103/analysis/spectrum/effective_mass_nopol_P200.png}

\source{P200 nopol effective mass；数据文件：effective\_mass\_nopol\_P200.png}\end{column}\end{columns}

  \source{\texttt{\detokenize{pyqcd/analysis/_analyse.py:253}}}
\end{frame}

\begin{frame}[t]{色散和连接 3pt 作为控制道，和非局域胶子 3pt 严格区分}
  \begin{columns}[T,totalwidth=\textwidth]
\begin{column}{0.49\textwidth}\includegraphics[width=0.96\linewidth,height=0.62\textheight,keepaspectratio]{/root/PyQCD/data/runs/tmd-ext-20261103/analysis/spectrum/dispersion_diagnostic.png}

\source{核子能量色散诊断；数据文件：dispersion\_diagnostic.png}\end{column}
\begin{column}{0.49\textwidth}\includegraphics[width=0.96\linewidth,height=0.62\textheight,keepaspectratio]{/root/PyQCD/data/runs/tmd-ext-20261103/analysis/spectrum/ratio_3pt_proton_P0.png}

\source{连接矢量流 3pt/2pt ratio 控制道；数据文件：ratio\_3pt\_proton\_P0.png}\end{column}\end{columns}

  \source{\texttt{\detokenize{pyqcd/analysis/_correlators.py:84}}}
\end{frame}

\begin{frame}[t]{方向与动量比较保留谱学平台的方向依赖}
  \includegraphics[width=0.88\linewidth,height=0.62\textheight,keepaspectratio]{/root/PyQCD/data/runs/tmd-ext-20261103/analysis/spectrum/direction_momentum_comparison.png}

\source{nopol/pol35 的五个动量与方向平台比较；数据文件：direction\_momentum\_comparison.png}

  \source{\texttt{\detokenize{pyqcd/testing/tmd/analyze_full_production.py:739}}}
\end{frame}

\begin{frame}[t]{有限横距 OPE 的离散量已完整落盘，三方向用于旋转对称性检查}
  \small\begin{tabularx}{\linewidth}{@{}l l X@{}}\toprule 对象 & shape & 作用 \\ \midrule 1D full-channel & $(2,4,6,13,7,72)$ & unpolarized/helicity 全通道 \\ 1D legacy & $(13,7,72)$ & 形状归一与交叉检查 \\ 2D full-channel & $(2,4,6,13,7,3,72)$ & 双分量 transverse staircase \\ flowed gauge & $(72,24,24,24,4,3,3)$ & $V_\tau$ 原始流场缓存 \\ \bottomrule\end{tabularx}

  \source{\texttt{\detokenize{pyqcd/pipeline/_tmd9.py:868}}}
\end{frame}

\begin{frame}[t]{三方向 ratio 的复相位与 helicity 通道被保留到统计层}
  \begin{columns}[T,totalwidth=\textwidth]
\begin{column}{0.32\textwidth}\includegraphics[width=0.98\linewidth,height=0.62\textheight,keepaspectratio]{/root/PyQCD/data/runs/tmd-ext-20261103/analysis/ratio_1d/ratio_v2_tau1_x_P002_ratio.png}

\source{x 方向 ratio；数据文件：ratio\_v2\_tau1\_x\_P002\_ratio.png}\end{column}
\begin{column}{0.32\textwidth}\includegraphics[width=0.98\linewidth,height=0.62\textheight,keepaspectratio]{/root/PyQCD/data/runs/tmd-ext-20261103/analysis/ratio_1d/ratio_v2_tau1_y_P020_ratio.png}

\source{y 方向 ratio；数据文件：ratio\_v2\_tau1\_y\_P020\_ratio.png}\end{column}
\begin{column}{0.32\textwidth}\includegraphics[width=0.98\linewidth,height=0.62\textheight,keepaspectratio]{/root/PyQCD/data/runs/tmd-ext-20261103/analysis/ratio_1d/ratio_v2_tau1_z_P200_ratio.png}

\source{z 方向 ratio；数据文件：ratio\_v2\_tau1\_z\_P200\_ratio.png}\end{column}\end{columns}

  \source{\texttt{\detokenize{pyqcd/analysis/_gradient_flow_ratio.py:980}}}
\end{frame}

\begin{frame}[t]{二维 staircase 顺序差异作为显式离散化系统误差报告}
  \includegraphics[width=0.88\linewidth,height=0.62\textheight,keepaspectratio]{/root/PyQCD/data/runs/tmd-ext-20261103/analysis/ratio_2d/staircase_sensitivity.png}

\source{order (0,1) 与 (1,0) ratio 比较；数据文件：staircase\_sensitivity.png}

\small 中位相对 L2：0.2203

  \source{\texttt{\detokenize{pyqcd/testing/tmd/analyze_full_production.py:1016}}}
\end{frame}

\begin{frame}[t]{二维 c0 同时保留 $b_x,b_y$ 依赖，避免把不同 staircase 强行合并}
  \begin{columns}[T,totalwidth=\textwidth]
\begin{column}{0.49\textwidth}\includegraphics[width=0.96\linewidth,height=0.62\textheight,keepaspectratio]{/root/PyQCD/data/runs/tmd-ext-20261103/analysis/plateaus_2d/c0_tau1_order01_heatmap.png}

\source{tau=1 2D c0 heatmap；数据文件：c0\_tau1\_order01\_heatmap.png}\end{column}
\begin{column}{0.49\textwidth}\includegraphics[width=0.96\linewidth,height=0.62\textheight,keepaspectratio]{/root/PyQCD/data/runs/tmd-ext-20261103/analysis/plateaus_2d/c0_tau3_order01_heatmap.png}

\source{tau=3 2D c0 heatmap；数据文件：c0\_tau3\_order01\_heatmap.png}\end{column}\end{columns}

  \source{\texttt{\detokenize{pyqcd/testing/tmd/analyze_full_production.py:1162}}}
\end{frame}

\begin{frame}[t]{流时 $\tau/a^2=1$ 的 c0 平台对窗口敏感，不能由单窗直接宣称物理矩阵元}
  \includegraphics[width=0.88\linewidth,height=0.62\textheight,keepaspectratio]{/root/PyQCD/data/runs/tmd-ext-20261103/analysis/plateaus/c0_tau1_z_P200_windows.png}

\source{tau=1 多窗口 c0 稳定性；数据文件：c0\_tau1\_z\_P200\_windows.png}

  \source{窗口结果：\nolinkurl{/root/PyQCD/data/runs/tmd-ext-20261103/analysis/plateaus/c0_tau1_z_P200.npz}}
\end{frame}

\begin{frame}[t]{流时 $\tau/a^2=2$ 的 c0 平台对窗口敏感，不能由单窗直接宣称物理矩阵元}
  \includegraphics[width=0.88\linewidth,height=0.62\textheight,keepaspectratio]{/root/PyQCD/data/runs/tmd-ext-20261103/analysis/plateaus/c0_tau2_z_P200_windows.png}

\source{tau=2 多窗口 c0 稳定性；数据文件：c0\_tau2\_z\_P200\_windows.png}

  \source{窗口结果：\nolinkurl{/root/PyQCD/data/runs/tmd-ext-20261103/analysis/plateaus/c0_tau2_z_P200.npz}}
\end{frame}

\begin{frame}[t]{流时 $\tau/a^2=3$ 的 c0 平台对窗口敏感，不能由单窗直接宣称物理矩阵元}
  \includegraphics[width=0.88\linewidth,height=0.62\textheight,keepaspectratio]{/root/PyQCD/data/runs/tmd-ext-20261103/analysis/plateaus/c0_tau3_z_P200_windows.png}

\source{tau=3 多窗口 c0 稳定性；数据文件：c0\_tau3\_z\_P200\_windows.png}

  \source{窗口结果：\nolinkurl{/root/PyQCD/data/runs/tmd-ext-20261103/analysis/plateaus/c0_tau3_z_P200.npz}}
\end{frame}

\begin{frame}[t]{Sumratio AIC 质量与 fallback 覆盖作为独立诊断层报告}
  \includegraphics[width=0.88\linewidth,height=0.62\textheight,keepaspectratio]{/root/PyQCD/data/runs/tmd-ext-20261103/analysis/plateaus/sumratio_quality.png}

\source{Sumratio quality-accepted fraction 总览；数据文件：sumratio\_quality.png}

  \source{\texttt{\detokenize{pyqcd/testing/tmd/analyze_full_production.py:1134}}}
\end{frame}

\begin{frame}[t]{流窗交集只保留可比较几何，跨流时差异采用配对 JK}
  \begin{columns}[T,totalwidth=\textwidth]
\begin{column}{0.49\textwidth}\includegraphics[width=0.96\linewidth,height=0.62\textheight,keepaspectratio]{/root/PyQCD/data/runs/tmd-ext-20261103/analysis/flow_windows/flow_window_coverage.png}

\source{流窗覆盖；数据文件：flow\_window\_coverage.png}\end{column}
\begin{column}{0.49\textwidth}\includegraphics[width=0.96\linewidth,height=0.62\textheight,keepaspectratio]{/root/PyQCD/data/runs/tmd-ext-20261103/analysis/flow_windows/flow_time_paired_difference.png}

\source{配对流时差；数据文件：flow\_time\_paired\_difference.png}\end{column}\end{columns}

  \source{共同窗：\{'z\_indices': [5, 6, 7, 8, 9, 10, 11, 12], 'b\_indices': [5, 6], 'n\_points': 16\}。}
\end{frame}

\begin{frame}[t]{格点端准 PDF 与准 TMD-PDF 仅报告诊断状态，未施加光锥匹配}
  \begin{columns}[T,totalwidth=\textwidth]
\begin{column}{0.49\textwidth}\includegraphics[width=0.96\linewidth,height=0.62\textheight,keepaspectratio]{/root/PyQCD/data/runs/tmd-ext-20261103/analysis/pdf/quasi_tau1_P200.png}

\source{tau=1, P200 准 PDF/TMD-PDF；数据文件：quasi\_tau1\_P200.png}\end{column}
\begin{column}{0.49\textwidth}\includegraphics[width=0.96\linewidth,height=0.62\textheight,keepaspectratio]{/root/PyQCD/data/runs/tmd-ext-20261103/analysis/pdf/quasi_tau3_P200.png}

\source{tau=3, P200 准 PDF/TMD-PDF；数据文件：quasi\_tau3\_P200.png}\end{column}\end{columns}

  \source{\texttt{\detokenize{pyqcd/renorm/_tmdextract.py:136}}}
\end{frame}

\begin{frame}[t]{重整化实空间矩阵元 $h_R(z,b)$ 保留 $b$ 依赖并区分流时}
  \begin{columns}[T,totalwidth=\textwidth]
\begin{column}{0.49\textwidth}\includegraphics[width=0.96\linewidth,height=0.62\textheight,keepaspectratio]{/root/PyQCD/data/runs/tmd-ext-20261103/analysis/pdf/hR_tau1_P200.png}

\source{tau=1, P200 hR(z,b)；数据文件：hR\_tau1\_P200.png}\end{column}
\begin{column}{0.49\textwidth}\includegraphics[width=0.96\linewidth,height=0.62\textheight,keepaspectratio]{/root/PyQCD/data/runs/tmd-ext-20261103/analysis/pdf/hR_tau3_P200.png}

\source{tau=3, P200 hR(z,b)；数据文件：hR\_tau3\_P200.png}\end{column}\end{columns}

  \source{\texttt{\detokenize{pyqcd/pipeline/_tmd9.py:1407}}}
\end{frame}

\begin{frame}[t]{P400 准 PDF/TMD-PDF 与 hR 作为高动量诊断单独展示}
  \begin{columns}[T,totalwidth=\textwidth]
\begin{column}{0.49\textwidth}\includegraphics[width=0.96\linewidth,height=0.62\textheight,keepaspectratio]{/root/PyQCD/data/runs/tmd-ext-20261103/analysis/pdf/quasi_tau3_P400.png}

\source{tau=3, P400 准 PDF/TMD-PDF；数据文件：quasi\_tau3\_P400.png}\end{column}
\begin{column}{0.49\textwidth}\includegraphics[width=0.96\linewidth,height=0.62\textheight,keepaspectratio]{/root/PyQCD/data/runs/tmd-ext-20261103/analysis/pdf/hR_tau3_P400.png}

\source{tau=3, P400 hR(z,b)；数据文件：hR\_tau3\_P400.png}\end{column}\end{columns}

  \source{\texttt{\detokenize{pyqcd/renorm/_tmdextract.py:32}}}
\end{frame}

\begin{frame}[t]{两动量 CS 核状态与信噪比门控作为格点诊断保存}
  \begin{columns}[T,totalwidth=\textwidth]
\begin{column}{0.49\textwidth}\includegraphics[width=0.96\linewidth,height=0.62\textheight,keepaspectratio]{/root/PyQCD/data/runs/tmd-ext-20261103/analysis/pdf/cs_kernel_tau1.png}

\source{tau=1 CS kernel；数据文件：cs\_kernel\_tau1.png}\end{column}
\begin{column}{0.49\textwidth}\includegraphics[width=0.96\linewidth,height=0.62\textheight,keepaspectratio]{/root/PyQCD/data/runs/tmd-ext-20261103/analysis/pdf/cs_kernel_tau3.png}

\source{tau=3 CS kernel；数据文件：cs\_kernel\_tau3.png}\end{column}\end{columns}

  \source{\texttt{\detokenize{pyqcd/renorm/_tmdextract.py:357}}}
\end{frame}

\begin{frame}[t]{误差账本把统计、平台、流时、staircase 与未实现层分开}
  \small\begin{tabularx}{\linewidth}{@{}l l X@{}}\toprule 来源 & 数值 & 解释 \\ \midrule 统计中位 & 0.22 & JK plateau 误差 \\ 平台窗漂移 & 1.215 & 多窗口极差 \\ 流时相对效应 & 0.8944 & 配对 JK 差值 \\ staircase 相对 L2 & 0.2203 & 路径顺序系统差 \\ soft/rapidity & 未测量 & PDF 不得物理解释 \\ matching/mixing & 未施加 & 保持格点端诊断 \\ \bottomrule\end{tabularx}

  \source{完整账本：\nolinkurl{/root/PyQCD/data/runs/tmd-ext-20261103/analysis/error_budget.json}}
\end{frame}

\begin{frame}[t]{误差账本的量级用于判断当前主导不确定性，而不是用于升级物理结论}
  \includegraphics[width=0.88\linewidth,height=0.62\textheight,keepaspectratio]{/root/PyQCD/data/runs/tmd-ext-20261103/analysis/error_budget.png}

\source{统计、窗口、流时与 staircase 误差量级；数据文件：error\_budget.png}

  \source{完整账本：\nolinkurl{/root/PyQCD/data/runs/tmd-ext-20261103/analysis/error_budget.json}}
\end{frame}

\begin{frame}[t]{当前结论：数值链闭合为 diagnostic-only，物理闭环明确失败关闭}
  \begin{columns}[T,totalwidth=\textwidth]
\begin{column}{0.48\textwidth}\begin{block}{可确认}13 组态真实数据、三流时、三方向与 2D OPE、ratio 和格点端 Fourier 均有可读产物。\end{block}\end{column}
\begin{column}{0.48\textwidth}\begin{block}{不可确认}不能将当前 xg 解读为物理胶子 TMD-PDF；流时/platform/matching 等系统误差仍主导。\end{block}\end{column}\end{columns}

  \source{\texttt{\detokenize{pyqcd/testing/tmd/analyze_full_production.py:213}}}
\end{frame}

\begin{frame}[t]{下一步按物理瓶颈排序：soft/rapidity、系统扫描、连续极限}
  \small\begin{tabularx}{\linewidth}{@{}l l X@{}}\toprule 任务 & 验收产物 & 通过条件 \\ \midrule soft/rapidity & 同几何软因子与误差 & 表示、流时、staple 完全匹配 \\ tau/staple 扫描 & 多尺度 hR 与外推带 & 共同窗内稳定且误差账本闭合 \\ matching & 显式胶子混合矩阵 & $\alpha_s\to0$ 核对角核 \\ 连续极限 & 多格距/体积联合外推 & 统计与系统误差分离 \\ \bottomrule\end{tabularx}

  \source{本任务按要求未执行非格点 QCD matching。}
\end{frame}

\begin{frame}[t]{附录：物理对象到代码入口的完整映射}
  \begin{tabularx}{\linewidth}{@{}l X@{}}
\toprule
物理/工程对象 & 代码证据 \\
\midrule
Wilson flow RK3 & \texttt{\detokenize{pyqcd/renorm/_gradient_flow.py:176}}\\
Clover/staple TMD OPE & \texttt{\detokenize{pyqcd/renorm/_tmd.py:366}}\\
canonical 1D/2D OPE & \texttt{\detokenize{pyqcd/pipeline/_tmd9.py:1114}}\\
disconnected ratio & \texttt{\detokenize{pyqcd/analysis/_gradient_flow_ratio.py:607}}\\
plateau matrix element & \texttt{\detokenize{pyqcd/analysis/_tmd_ratio.py:986}}\\
quasi-TMD Fourier & \texttt{\detokenize{pyqcd/renorm/_tmdextract.py:32}}\\
lattice CS kernel & \texttt{\detokenize{pyqcd/renorm/_tmdextract.py:357}}\\
production driver & \texttt{\detokenize{pyqcd/testing/tmd/run_full_production.py:693}}\\
analysis driver & \texttt{\detokenize{pyqcd/testing/tmd/analyze_full_production.py:1741}}\\
\bottomrule
\end{tabularx}

  \source{所有入口均位于 PyQCD 生产包；测试驱动只负责编排和持久化。}
\end{frame}

\begin{frame}[t]{附录：产物与复现入口}
  \small\begin{tabularx}{\linewidth}{@{}l X@{}}\toprule 对象 & 路径/命令 \\ \midrule 运行根目录 & \nolinkurl{/root/PyQCD/data/runs/tmd-ext-20261103} \\ 分析汇总 & \nolinkurl{/root/PyQCD/data/runs/tmd-ext-20261103/analysis/analysis_summary.json} \\ 产物清单 & \nolinkurl{/root/PyQCD/data/runs/tmd-ext-20261103/analysis/artifact_manifest.json} \\ 生产入口 & \nolinkurl{/root/PyQCD/pyqcd/testing/tmd/run_full_production.py} \\ 分析入口 & \nolinkurl{/root/PyQCD/pyqcd/testing/tmd/analyze_full_production.py} \\ \bottomrule\end{tabularx}

  \source{复现时先读取 workers checkpoint，再调用两个入口的相同请求参数。}
\end{frame}

\end{document}
